
\documentclass[11pt,a4paper]{article}

% Shared macros for the theory notes.
% Include with % Shared macros for the theory notes.
% Include with % Shared macros for the theory notes.
% Include with \input{macros} after \begin{document} preamble packages.

\newcommand{\Om}{\Omega}
\newcommand{\dOm}{\partial\Omega}
\newcommand{\dx}{\,\mathrm{d}x}
\newcommand{\ds}{\,\mathrm{d}s}

\newcommand{\Ltwo}{L^2(\Om)}
\newcommand{\Ltwozero}{L^2_0(\Om)}
\newcommand{\Hone}{H^1(\Om)}
\newcommand{\Honezero}{H^1_0(\Om)}

\newcommand{\normL}[1]{\left\lVert #1 \right\rVert_{L^2(\Om)}}
\newcommand{\normH}[1]{\left\lVert #1 \right\rVert_{H^1(\Om)}}
\newcommand{\ip}[2]{\left( #1 , #2 \right)}

\newcommand{\grad}{\nabla}
\newcommand{\lap}{\Delta}
\newcommand{\dn}[1]{\frac{\partial #1}{\partial n}}

\newtheorem{theorem}{Theorem}[section]
\newtheorem{lemma}[theorem]{Lemma}
\newtheorem{proposition}[theorem]{Proposition}
\theoremstyle{definition}
\newtheorem{definition}[theorem]{Definition}
\newtheorem{remark}[theorem]{Remark}
 after \begin{document} preamble packages.

\newcommand{\Om}{\Omega}
\newcommand{\dOm}{\partial\Omega}
\newcommand{\dx}{\,\mathrm{d}x}
\newcommand{\ds}{\,\mathrm{d}s}

\newcommand{\Ltwo}{L^2(\Om)}
\newcommand{\Ltwozero}{L^2_0(\Om)}
\newcommand{\Hone}{H^1(\Om)}
\newcommand{\Honezero}{H^1_0(\Om)}

\newcommand{\normL}[1]{\left\lVert #1 \right\rVert_{L^2(\Om)}}
\newcommand{\normH}[1]{\left\lVert #1 \right\rVert_{H^1(\Om)}}
\newcommand{\ip}[2]{\left( #1 , #2 \right)}

\newcommand{\grad}{\nabla}
\newcommand{\lap}{\Delta}
\newcommand{\dn}[1]{\frac{\partial #1}{\partial n}}

\newtheorem{theorem}{Theorem}[section]
\newtheorem{lemma}[theorem]{Lemma}
\newtheorem{proposition}[theorem]{Proposition}
\theoremstyle{definition}
\newtheorem{definition}[theorem]{Definition}
\newtheorem{remark}[theorem]{Remark}
 after \begin{document} preamble packages.

\newcommand{\Om}{\Omega}
\newcommand{\dOm}{\partial\Omega}
\newcommand{\dx}{\,\mathrm{d}x}
\newcommand{\ds}{\,\mathrm{d}s}

\newcommand{\Ltwo}{L^2(\Om)}
\newcommand{\Ltwozero}{L^2_0(\Om)}
\newcommand{\Hone}{H^1(\Om)}
\newcommand{\Honezero}{H^1_0(\Om)}

\newcommand{\normL}[1]{\left\lVert #1 \right\rVert_{L^2(\Om)}}
\newcommand{\normH}[1]{\left\lVert #1 \right\rVert_{H^1(\Om)}}
\newcommand{\ip}[2]{\left( #1 , #2 \right)}

\newcommand{\grad}{\nabla}
\newcommand{\lap}{\Delta}
\newcommand{\dn}[1]{\frac{\partial #1}{\partial n}}

\newtheorem{theorem}{Theorem}[section]
\newtheorem{lemma}[theorem]{Lemma}
\newtheorem{proposition}[theorem]{Proposition}
\theoremstyle{definition}
\newtheorem{definition}[theorem]{Definition}
\newtheorem{remark}[theorem]{Remark}


\title{Weak Formulations of Scalar Elliptic Problems:\\
Boundedness, Coercivity, and How the Function Space Is Chosen}
\author{M. Javed Bukhari}
\date{\today}

\begin{document}
\maketitle

\begin{abstract}
These notes work through three scalar elliptic boundary value problems. For each
one I derive the weak formulation and then prove that the resulting bilinear form
is bounded and coercive. Two things are worth watching. The first is that the
function space is not chosen in advance --- it falls out of the weak formulation
once we ask what is needed to make the integrals finite. The second is that
boundedness and coercivity point in opposite directions and need different tools:
boundedness is an upper bound on $a(u,v)$ and comes from Cauchy--Schwarz, while
coercivity is a lower bound on $a(u,u)$ and comes from setting $v=u$ and then
finding the weakest coefficient. The third problem is the microrotation equation
of the micropolar system with the coupling switched off, so the last result here
is used directly later on.
\end{abstract}

\tableofcontents

%==============================================================================
\section{Notation and assumptions}
%==============================================================================

Throughout, $\Om \subset \R^{n}$ is a bounded Lipschitz domain with boundary
$\dOm$. All derivatives are weak derivatives. I write
\[
  \normL{u} = \left( \int_\Om \abs{u}^{2} \dx \right)^{1/2},
  \qquad
  \normH{u}^{2} = \normL{u}^{2} + \normL{\grad u}^{2},
\]
and $\Honezero$ for the closure of $\Cinfc$ in the $\Hone$ norm. When $n=1$ and
$\Om=(0,1)$ I write $u'$ instead of $\grad u$.

Two results are used over and over.

\begin{lemma}[Cauchy--Schwarz]\label{lem:cs}
For $u,v \in \Ltwo$,
\[
  \abs{\int_\Om uv \dx} \;\le\; \normL{u}\,\normL{v}.
\]
\end{lemma}

\begin{lemma}[Poincar\'e--Friedrichs]\label{lem:poincare}
There is a constant $C_\Om > 0$, depending only on $\Om$, with
\[
  \normL{u} \;\le\; C_\Om \normL{\grad u}
  \qquad \text{for every } u \in \Honezero .
\]
\end{lemma}

\begin{remark}\label{rem:poincare-fails}
Lemma~\ref{lem:poincare} is false on $\Hone$. Take $u \equiv 1$. Then
$\grad u = 0$, so the right-hand side is zero, but
$\normL{u} = \abs{\Om}^{1/2} \neq 0$. What makes the inequality work on
$\Honezero$ is that the boundary condition rules out non-zero constants. So the
boundary condition is not a side detail. It is the hypothesis the whole existence
theory rests on, and I come back to this point at the end of
Section~\ref{sec:reaction}.
\end{remark}

\begin{remark}
One consequence of Lemma~\ref{lem:poincare} is that $\normL{\grad \cdot}$ is
itself a norm on $\Honezero$, equivalent to $\normH{\cdot}$. Many papers use it
without saying so, which is confusing the first time you meet it.
\end{remark}

%==============================================================================
\section{The one-dimensional model problem}
%==============================================================================

\subsection{From the strong form to the weak form}

Take the two-point boundary value problem
\begin{equation}\label{eq:strongA}
  -u'' = f \quad \text{on } (0,1), \qquad u(0) = u(1) = 0 .
\end{equation}

Let $v$ be a test function with $v(0)=v(1)=0$. Multiply \eqref{eq:strongA} by $v$
and integrate over $(0,1)$:
\[
  -\int_0^1 u'' v \dx = \int_0^1 f v \dx .
\]
Integrate the left-hand side by parts once:
\[
  -\int_0^1 u'' v \dx
  = -\Big[\, u' v \,\Big]_0^1 + \int_0^1 u' v' \dx
  = \int_0^1 u' v' \dx ,
\]
where the boundary term drops out because $v(0)=v(1)=0$. It is worth being
careful about which function is doing that work: it is $v$, not $u'$. There is no
reason for $u'$ to vanish at the endpoints, and in general it does not.

So the weak problem is: find $u \in H^1_0(0,1)$ with
\begin{equation}\label{eq:weakA}
  a(u,v) \;:=\; \int_0^1 u' v' \dx
  \;=\; \int_0^1 f v \dx \;=:\; F(v)
  \qquad \text{for all } v \in H^1_0(0,1).
\end{equation}

The trade we made is worth naming. The classical form needs two derivatives of
$u$ and none of $v$; the weak form needs one of each. We gave away regularity of
$v$, which costs nothing because we choose $v$ ourselves and can always choose it
smooth, in order to buy back regularity of $u$, which is what we actually care
about.

\subsection{Why the space is $H^1_0$}
\label{sec:space-choice}

The classical formulation \eqref{eq:strongA} asks for $u \in C^{2}$. That demand
comes from the way we wrote the equation down, not from the problem being
modelled, and it throws away solutions that exist in reality. If $f$ has a jump
--- a bar under a point load, or a change of material halfway along --- then no
$C^{2}$ function satisfies \eqref{eq:strongA}. The bar still bends. So a theory
built on $C^{2}$ would tell us there is no solution to a problem whose solution
we can see.

The weak form \eqref{eq:weakA} only involves $u'$, so the right space is found by
asking what is needed for its integrals to make sense. By Lemma~\ref{lem:cs},
$\int_0^1 u'v'$ is finite as soon as $u', v' \in L^{2}(0,1)$, and we needed
$v(0)=v(1)=0$ to kill the boundary term. That gives
\[
  H^1_0(0,1) = \{\, u \in L^{2}(0,1) : u' \in L^{2}(0,1),\ u(0)=u(1)=0 \,\} .
\]
Note the direction of that argument. The space was built out of the requirement;
it was not picked first and then imposed.

The choice is squeezed from both sides. Go bigger --- drop the condition
$u' \in L^{2}$ and work in $L^{2}(0,1)$ alone --- and $\int_0^1 u'v'$ may be
infinite, so the equation stops meaning anything. Go smaller --- $C^{2}$, or
$H^{2} \cap H^1_0$ --- and we are assuming more regularity than the formulation
ever asked for, which throws out the point-load solution above. $H^1_0$ is the
largest space on which \eqref{eq:weakA} still makes sense, and that is what makes
it the right one.

\begin{remark}
In one dimension $H^{1}(0,1) \hookrightarrow C[0,1]$, so writing
$u(0)=u(1)=0$ is fine. In two dimensions this needs more care, because $\dOm$
has measure zero and $\Hone$ functions are only defined up to null sets. See
Remark~\ref{rem:trace}.
\end{remark}

\subsection{Boundedness}

\begin{proposition}\label{prop:boundA}
For all $u,v \in H^1_0(0,1)$ we have $\abs{a(u,v)} \le \normH{u}\,\normH{v}$.
\end{proposition}

\begin{proof}
Cauchy--Schwarz (Lemma~\ref{lem:cs}) applied to $u'$ and $v'$ gives
\begin{equation}\label{eq:boundA1}
  \abs{a(u,v)} = \abs{\int_0^1 u'v' \dx} \;\le\; \normL{u'}\,\normL{v'} .
\end{equation}
Now $\normH{u}^{2} = \normL{u}^{2} + \normL{u'}^{2}$ and $\normL{u}^{2} \ge 0$,
so $\normL{u'}^{2} \le \normH{u}^{2}$, i.e.
\begin{equation}\label{eq:boundA2}
  \normL{u'} \le \normH{u},
  \qquad \text{and in the same way} \qquad
  \normL{v'} \le \normH{v}.
\end{equation}
Putting \eqref{eq:boundA1} and \eqref{eq:boundA2} together finishes it, with
constant $C=1$.
\end{proof}

\begin{remark}
Poincar\'e does not appear here, and it could not have. Boundedness needs the
seminorm controlled by the full norm, $\normL{u'} \le \normH{u}$, which is free.
Poincar\'e gives control in the opposite direction, $\normH{u} \lesssim
\normL{u'}$, which is not free and is what coercivity needs. Two opposite
directions, two different jobs. Mixing them up is the easiest mistake to make
here, and I made it the first time.
\end{remark}

\subsection{Coercivity}

\begin{proposition}\label{prop:coercA}
For all $u \in H^1_0(0,1)$,
\[
  a(u,u) \;\ge\; \alpha \normH{u}^{2},
  \qquad \alpha = \frac{1}{C_\Om^{2}+1} .
\]
\end{proposition}

\begin{proof}
Put $v=u$. This gives an identity, not an estimate:
\begin{equation}\label{eq:coercA1}
  a(u,u) = \int_0^1 (u')^{2} \dx = \normL{u'}^{2} .
\end{equation}
There is nothing to bound at this step --- $\int_0^1 (u')^{2}$ \emph{is}
$\normL{u'}^{2}$, straight from the definition of the norm. This is why setting
$v=u$ is the whole trick: Cauchy--Schwarz exists to handle two different
functions, and with the same function twice there is nothing left to estimate.

By Lemma~\ref{lem:poincare}, $\normL{u}^{2} \le C_\Om^{2}\normL{u'}^{2}$, so
\begin{equation}\label{eq:coercA2}
  \normH{u}^{2} = \normL{u}^{2} + \normL{u'}^{2}
  \;\le\; C_\Om^{2}\normL{u'}^{2} + \normL{u'}^{2}
  = \left( C_\Om^{2}+1 \right)\normL{u'}^{2} .
\end{equation}
Rearranging \eqref{eq:coercA2},
\begin{equation}\label{eq:coercA3}
  \normL{u'}^{2} \;\ge\; \frac{1}{C_\Om^{2}+1}\,\normH{u}^{2} ,
\end{equation}
and combining this with \eqref{eq:coercA1} gives the result.
\end{proof}

\begin{remark}[Sanity check]
The constant is $\alpha = (C_\Om^{2}+1)^{-1}$, so $\alpha$ gets smaller as
$C_\Om$ gets bigger, which means coercivity gets weaker on domains with a large
Poincar\'e constant. Since $C_\Om$ grows with the size of $\Om$, this is saying
something I find intuitive: on a long interval the two endpoint conditions do not
constrain the function much, because a function with a small derivative still has
room to wander a long way from zero in the middle. On a short interval the same
two conditions clamp it. So $\alpha$ is measuring how hard the boundary condition
is working, which lines up with Remark~\ref{rem:poincare-fails}.
\end{remark}

\subsection{Stability of the solution}

\begin{proposition}\label{prop:stabA}
Let $f \in L^{2}(0,1)$ and let $u \in H^1_0(0,1)$ solve \eqref{eq:weakA}. Then
\[
  \normH{u} \;\le\; \frac{1}{\alpha}\normL{f},
  \qquad \alpha = \frac{1}{C_\Om^{2}+1} .
\]
\end{proposition}

\begin{proof}
Take $v=u$ in \eqref{eq:weakA}, so $a(u,u) = \int_0^1 fu \dx$. Bound the left
side from below with Proposition~\ref{prop:coercA} and the right side from above
with Lemma~\ref{lem:cs}:
\[
  \alpha\normH{u}^{2} \;\le\; a(u,u) = \int_0^1 fu \dx
  \;\le\; \normL{f}\normL{u} \;\le\; \normL{f}\normH{u} .
\]
If $u=0$ there is nothing to prove. Otherwise divide by $\normH{u} > 0$.
\end{proof}

\begin{remark}
This says: bounded data gives a bounded solution, with a constant we can write
down. That is stability, and it is the continuous version of what we want from a
numerical scheme --- that it should not amplify perturbations in the data. Worth
noticing that it came from coercivity and not from boundedness. Coercivity is the
property that does the real work; boundedness is a technical hypothesis that is
almost always easy.
\end{remark}

%==============================================================================
\section{The Dirichlet problem in two dimensions}
%==============================================================================

\subsection{From the strong form to the weak form}

Now let $\Om \subset \R^{2}$ and take
\begin{equation}\label{eq:strongB}
  -\lap u = f \ \text{ in } \Om, \qquad u = 0 \ \text{ on } \dOm .
\end{equation}

The tool is Green's first identity: for $u \in \Htwo$ and $v \in \Hone$,
\begin{equation}\label{eq:green}
  \int_\Om (\lap u)\,v \dx
  = \int_{\dOm} \dn{u}\,v \ds - \int_\Om \grad u \cdot \grad v \dx ,
\end{equation}
with $\dn{u} = \grad u \cdot \bn$ and $\bn$ the outward unit normal. This is
integration by parts in $n$ dimensions, and it lines up term by term with the
one-dimensional case: $\int_\Om (\lap u)v$ against $\int_0^1 u''v$, the boundary
integral against the evaluation $[u'v]_0^1$, and
$\int_\Om \grad u \cdot \grad v$ against $\int_0^1 u'v'$. In 1D the ``boundary''
is two points and the outward normal derivative is $+u'(1)$ and $-u'(0)$, which
is exactly what the square bracket computes.

Multiply \eqref{eq:strongB} by $v \in \Honezero$, integrate over $\Om$, and apply
\eqref{eq:green}:
\[
  -\int_\Om (\lap u)v \dx
  = -\int_{\dOm} \dn{u}\,v \ds + \int_\Om \grad u \cdot \grad v \dx
  = \int_\Om \grad u \cdot \grad v \dx ,
\]
the boundary integral vanishing because $v$ has zero trace. Again it is $v$ that
does this, not $\dn{u}$.

The weak problem: find $u \in \Honezero$ with
\begin{equation}\label{eq:weakB}
  a(u,v) := \int_\Om \grad u \cdot \grad v \dx
  = \int_\Om fv \dx
  \qquad \text{for all } v \in \Honezero ,
\end{equation}
where $\grad u \cdot \grad v
= \partial_x u\,\partial_x v + \partial_y u\,\partial_y v$. In one dimension that
sum has a single term and reads $u'v'$; in two dimensions it has two. That is the
only change.

\begin{remark}\label{rem:trace}
``$v=0$ on $\dOm$'' cannot be read pointwise here. The boundary has measure zero
in $\R^{2}$, and elements of $\Hone$ are equivalence classes modulo null sets, so
the value of $v$ at any particular point is not well defined. The statement is
made precise by the trace theorem: the restriction map
$v \mapsto v\vert_{\dOm}$, defined on $C(\overline{\Om}) \cap \Hone$, extends to a
bounded operator on $\Hone$, and $\Honezero$ is exactly the kernel of that
operator. I am citing this, not proving it.
\end{remark}

\begin{remark}
Identity \eqref{eq:green} needs $u \in \Htwo$ so that $\lap u$ is a function. That
regularity is used only to \emph{get to} \eqref{eq:weakB} from
\eqref{eq:strongB}. The weak problem itself only needs $u \in \Honezero$, and the
analysis below is done there. Deriving under classical regularity and then
dropping it is the standard order of business, and it is easy to be confused by
it the first time.
\end{remark}

\subsection{Boundedness and coercivity}

\begin{proposition}\label{prop:boundB}
For all $u,v \in \Honezero$, $\abs{a(u,v)} \le \normH{u}\normH{v}$.
\end{proposition}

\begin{proof}
Cauchy--Schwarz is used twice: once pointwise on the vectors $\grad u$ and
$\grad v$, then once on the integral via Lemma~\ref{lem:cs}. That gives
\[
  \abs{a(u,v)} = \abs{\int_\Om \grad u \cdot \grad v \dx}
  \le \int_\Om \abs{\grad u}\abs{\grad v} \dx
  \le \normL{\grad u}\normL{\grad v} .
\]
Then $\normL{\grad u} \le \normH{u}$ and $\normL{\grad v} \le \normH{v}$, exactly
as in \eqref{eq:boundA2}.
\end{proof}

\begin{proposition}\label{prop:coercB}
For all $u \in \Honezero$, $a(u,u) \ge \alpha\normH{u}^{2}$ with
$\alpha = (C_\Om^{2}+1)^{-1}$.
\end{proposition}

\begin{proof}
Setting $v=u$ gives
$a(u,u) = \int_\Om \abs{\grad u}^{2}\dx = \normL{\grad u}^{2}$. Poincar\'e holds
on $\Honezero$ in any dimension, so
\[
  \normH{u}^{2} = \normL{u}^{2} + \normL{\grad u}^{2}
  \le \left( C_\Om^{2}+1 \right)\normL{\grad u}^{2},
\]
and the rest is as in Proposition~\ref{prop:coercA}.
\end{proof}

\begin{remark}
These two proofs are the one-dimensional ones with $u'$ replaced by $\grad u$,
and the constants come out the same. The dimension could have caused trouble in
two places --- the integration by parts, and the meaning of the boundary
condition --- and both were dealt with by quoting a theorem (Green's identity,
and the trace theorem). Neither changed the shape of the argument. The point of
this section is that nothing new happens, which is itself worth knowing: the
machinery is dimension-independent.
\end{remark}

%==============================================================================
\section{The reaction--diffusion problem}
\label{sec:reaction}
%==============================================================================

\subsection{From the strong form to the weak form}

Let $c, \kappa > 0$ and take
\begin{equation}\label{eq:strongC}
  -c\,\lap u + \kappa u = f \ \text{ in } \Om, \qquad u = 0 \ \text{ on } \dOm .
\end{equation}

Multiply by $v \in \Honezero$ and integrate:
\[
  -c\int_\Om (\lap u)v \dx + \kappa\int_\Om uv \dx = \int_\Om fv \dx .
\]
Apply \eqref{eq:green} to the first term; its boundary contribution vanishes for
the same reason as before. The second term is left alone --- there is no
derivative in it, so there is nothing to integrate by parts. This gives: find
$u \in \Honezero$ with
\begin{equation}\label{eq:weakC}
  a(u,v) := c\int_\Om \grad u \cdot \grad v \dx + \kappa\int_\Om uv \dx
  = \int_\Om fv \dx
  \qquad \text{for all } v \in \Honezero .
\end{equation}

\subsection{Coercivity, this time without Poincar\'e}

\begin{proposition}\label{prop:coercC}
For all $u \in \Honezero$,
\[
  a(u,u) \;\ge\; \min(c,\kappa)\,\normH{u}^{2} .
\]
\end{proposition}

\begin{proof}
Set $v=u$ in \eqref{eq:weakC}. As before this is an identity:
\[
  a(u,u) = c\int_\Om \abs{\grad u}^{2}\dx + \kappa\int_\Om u^{2}\dx
  = c\,\normL{\grad u}^{2} + \kappa\,\normL{u}^{2} .
\]
Both pieces of $\normH{u}^{2}$ are already sitting there. The only difference is
the coefficient in front of each. Replacing both coefficients by the smaller of
the two,
\[
  a(u,u) \;\ge\; \min(c,\kappa)\left( \normL{\grad u}^{2} + \normL{u}^{2} \right)
  = \min(c,\kappa)\,\normH{u}^{2}. \qedhere
\]
\end{proof}

\begin{remark}
Poincar\'e is never used. The zeroth-order term hands us $\normL{u}^{2}$ for
free, and that is precisely the piece Poincar\'e had to manufacture in
Propositions~\ref{prop:coercA} and \ref{prop:coercB}.
\end{remark}

\begin{remark}[Sanity check]
As $\kappa \to 0$ the constant $\min(c,\kappa)$ goes to zero and this proof
collapses. That is not a flaw in the proof --- it is the proof telling us which
regime it belongs to. In that limit \eqref{eq:strongC} turns back into
\eqref{eq:strongB}, which is still coercive, but only via Poincar\'e and
therefore only via the boundary condition. So the two proofs are identifying two
different mechanisms for controlling $\normL{u}^{2}$. When $\kappa>0$ the
reaction term controls it directly and the boundary condition is not needed for
that job. When $\kappa=0$ the boundary condition is all there is, and it acts
through $C_\Om$. Something has to control $\normL{u}^{2}$; the question is only
what.
\end{remark}

%==============================================================================
\section{Summary}
%==============================================================================

\begin{center}
\begin{tabular}{@{}llll@{}}
\toprule
Problem & Space & Coercivity constant & What supplies coercivity \\
\midrule
$-u''=f$ & $H^1_0(0,1)$ & $(C_\Om^{2}+1)^{-1}$ & Poincar\'e \\
$-\lap u = f$ & $\Honezero$ & $(C_\Om^{2}+1)^{-1}$ & Poincar\'e \\
$-c\lap u + \kappa u = f$ & $\Honezero$ & $\min(c,\kappa)$ & the reaction term \\
\bottomrule
\end{tabular}
\end{center}

\medskip

Two things to take away.

First, in every case the space came out of the weak formulation rather than being
chosen beforehand. The recipe was always the same: write the weak form, look at
the integrals in it, and ask what the weakest condition on $u$ is that keeps them
finite. Section~\ref{sec:space-choice} then argued that this space is optimal ---
anything larger and the formulation stops making sense, anything smaller and we
lose solutions that genuinely exist.

Second, boundedness and coercivity are opposite kinds of statement and need
different tools. Boundedness bounds $a(u,v)$ from above, with two independent
arguments, and follows from Cauchy--Schwarz plus the free estimate
$\normL{\grad u} \le \normH{u}$. Coercivity bounds $a(u,u)$ from below, with the
same argument twice; setting $v=u$ produces an identity, and after that the only
question is what bounds $\normH{u}^{2}$ in terms of what the identity gave us.
Whether Poincar\'e is needed for that depends entirely on whether the equation
has a zeroth-order term. Once these two are in hand, Lax--Milgram gives existence
and uniqueness, which is the point of the whole exercise.

%==============================================================================
\section*{Where this connects to the micropolar problem}
\addcontentsline{toc}{section}{Where this connects to the micropolar problem}
%==============================================================================

Problem \eqref{eq:strongC} is the microrotation equation of the two-dimensional
stationary linear micropolar system,
\[
  -(c_a+c_d)\,\lap w + 4\nu_r\,w = 2\nu_r \curl \bu + g ,
\]
with the coupling term $2\nu_r \curl \bu$ switched off, so that $c = c_a+c_d$ and
$\kappa = 4\nu_r$. Proposition~\ref{prop:coercC} therefore identifies the term
that will have to absorb the indefinite coupling contribution when the full
system is analysed.

It also shows that the vortex viscosity $\nu_r$ enters the problem twice, and in
competing ways: it supplies the zeroth-order term that helps coercivity, and it
multiplies the coupling term that hurts it. Sorting out that competition ---
which requires Young's inequality with a free parameter, chosen so that the
coupling is absorbed without eating the terms that were controlling things ---
is the subject of the next set of notes.

\end{document}
